\documentclass[a4paper, 10pt]{scrartcl}


\usepackage{fontspec}
\newcommand{\mainFont}{NotoSerif}
\setmainfont{\mainFont}[
  Path           = ./fonts/\mainFont/static/,
  Extension      = .ttf,
  UprightFont    = *-Regular,
  BoldFont       = *-Bold,
  ItalicFont     = *-Italic,
  BoldItalicFont = *-BoldItalic,
]

\usepackage{xcolor}
\usepackage{tikz}
%\usepackage{lmodern} % Allows arbitrary font scaling without size limits
\usepackage{graphicx}
\usepackage{fontawesome5} % Icons: \faEnvelope, \faPhone, \faGithub, \faLinkedin, \faMapMarker, ...
\usepackage{hyperref} % Clickable mailto:/tel:/http links

\usetikzlibrary{calc} % Required for efficient coordinate calculations

\definecolor{background1}{HTML}{00171f}
\definecolor{background2}{HTML}{00335a}
\definecolor{col3}{HTML}{007da9}
\definecolor{picBorderCol}{HTML}{ffffff}
\definecolor{linkCol}{HTML}{490a83}

% Make links colored instead of the default red/boxed look
\hypersetup{
	colorlinks=true,
	urlcolor=linkCol,
	linkcolor=linkCol,
}

% options for text: \small, \normalsize, \large, \Large, \LARGE, \huge, \Huge

\pagestyle{empty}
\begin{document}

\hyphenpenalty=10000
\exhyphenpenalty=10000

\begin{tikzpicture}[remember picture, overlay, shift={(current page.north west)}]
	% (0,0) is now the top-left corner.




	% --- background ------------------------------------
	\pgfmathsetmacro{\splitpointX}{4}
	\pgfmathsetmacro{\splitpointY}{-2.8}
	\pgfmathsetmacro{\cornerpointUpX}{7.5}
	\pgfmathsetmacro{\cornerpointUpY}{-4.5}

	%\pgfmathsetmacro{\cornerpointDownX}{\cornerpointUpX}
	%\pgfmathsetmacro{\cornerpointDownY}{-7.8}

	\pgfmathsetmacro{\lastpointX}{4}
	\pgfmathsetmacro{\lastpointY}{-6.2}

	\fill[col3]
        (-1,1) --
        (-1,\lastpointY) --
        % Bézier S-curve connecting the bottom level to the top level
        (\lastpointX,\lastpointY) .. controls (\lastpointX + 2, \lastpointY) and (\cornerpointUpX - 2, \cornerpointUpY) .. (\cornerpointUpX,\cornerpointUpY) --
        (22,\cornerpointUpY) --
        (22,1) --
        cycle;


	%\fill[background2]
        %(-1,1) --
        %(-1,\splitpointY) --
        %% Bézier S-curve connecting the bottom level to the top level
        %(\splitpointX,\splitpointY) .. controls (\splitpointX + 2, \splitpointY) and (\cornerpointDownX - 2, \cornerpointDownY) .. (\cornerpointDownX,\cornerpointDownY) --
        %(22,\cornerpointDownY) --
        %(22,1) --
        %cycle;

	\fill[background1]
        (-1,1) --
        (-1,\splitpointY) --
        % Bézier S-curve connecting the bottom level to the top level
        (\splitpointX,\splitpointY) .. controls (\splitpointX + 2, \splitpointY) and (\cornerpointUpX - 2, \cornerpointUpY) .. (\cornerpointUpX,\cornerpointUpY) --
        (22,\cornerpointUpY) --
        (22,1) --
        cycle;













	% --- title name ------------------------------------
	\node[anchor=north west, text=white, font=\fontsize{42}{50}\selectfont\bfseries] at (7, -2.0)
	{
		Lucas Furer
	};







	% --- Contact section ------------------------------------
	% \faXXX gives the icon, \href{...}{...} makes it clickable.
	% mailto: opens the default mail client with the address
	% pre-filled; tel: does the same for a phone/dialer app.
	% On both PC and phone PDF readers you can usually also
	% right-click (or long-press) the link to copy it directly.
	\node[anchor=north west, align=left] at (0.5, -2.92)
	{
		\renewcommand{\arraystretch}{1.4}
		\begingroup
		\hypersetup{urlcolor=black}
		\begin{tabular}{@{}l l@{}}
			\faLinkedin   & \href{https://www.linkedin.com/in/lucas-furer-22a38a323/}{lucas-furer} \\
			\faGithub     & \href{https://github.com/LucasFurer}{LucasFurer} \\
			\faEnvelope   & \href{mailto:lucas.furer@gmail.com}{lucas.furer@gmail.com} \\
			\faPhone      & \href{tel:+31652772875}{+31 6 5277 2875} \\
			\faMapMarker  & Delft, NL \\
		\end{tabular}
		\endgroup
	};








	% --- section variables ------------------------------------
	\coordinate (sectStart) at ($(3.0,-7.2)$);


	\pgfmathsetmacro{\horGapSmall}{0.15} % space between text and the line
	\pgfmathsetmacro{\horGapLarge}{16.0} % space between line and dates
	\pgfmathsetmacro{\textWidth}{15.0} % width of the text

	\pgfmathsetmacro{\vertTextGap}{0.05} % spacing between text
	\pgfmathsetmacro{\vertSubGap}{0.1} % spacing between sub headers
	\pgfmathsetmacro{\vertSectGap}{0.3} % spacing between sections






	% --- text ------------------------------------
	\coordinate (eduStart) at ($(sectStart)$);


	%\node[anchor=north east, align=left, font=\Large] 
	%	at ($(eduStart) + (-\horGapSmall,0)$)
	%{
	%	\textbf{Education}
	%};


	\node[anchor=north west, align=left, text width=\textWidth cm] (eduEntryA) 
		at ($(eduStart) + (\horGapSmall,0)$)
	{
Beste selectiecommissie,\\
\vspace{1em}
Graag solliciteer ik naar het IT Young Professional Programma van Defensie. Recentelijk heb ik mijn master Computer Science aan de TU Delft afgerond, met een specialisatie in computergraphics. Deze functie spreekt mij aan omdat ik mijn technische kennis graag wil inzetten voor maatschappelijk relevant werk.\\
\vspace{1em}
Wat mij in het bijzonder aanspreekt in dit programma is de combinatie van werken aan concrete IT-uitdagingen en de ruimte om mij verder te ontwikkelen. Het Young Professional Programma biedt mij de mogelijkheid om ervaring op te doen, mijn sterke punten verder te ontdekken en mij geleidelijk te specialiseren in een richting die goed bij mij past.\\
\vspace{1em}
Tijdens mijn masterscriptie implementeerde ik in C++ een fast multipole N-body solver en integreerde deze in t-SNE. Daarnaast ontwikkelde ik een visualisatielaag met OpenGL, waarmee de resultaten direct konden worden geanalyseerd. Voor mijn bachelorscriptie onderzocht en implementeerde ik eveneens in C++ hoe de verdeling van lichtstraalreflecties verandert bij het gladder maken van microstructuren. Ook ontwikkelde ik zelfstandig een Blender addon die terrein genereert met gesimuleerde hydraulische erosie. Daarnaast werkte ik aan een Agile-project waarin ik de rol van Scrum Master vervulde en ervaring opdeed met GitLab CI/CD. Deze ervaringen hebben mij een sterke basis gegeven in zowel technisch uitdagend C++ werk als het ontwikkelen van complete softwareoplossingen.\\
\vspace{1em}
Naast mijn technische vaardigheden breng ik ook ervaring met samenwerken en het nemen van verantwoordelijkheid mee. Tijdens mijn stage bij Feedback Analytics gaf ik leiding aan een team van vijf personen en werkte ik nauw samen met de klant om de wensen en eisen in kaart te brengen. Deze ervaring heeft mij geleerd om helder te communiceren, actief mee te denken over oplossingen en verantwoordelijkheid te nemen voor de voortgang van een project.\\
\vspace{1em}
Ik zie mezelf als een nieuwsgierige en ambitieuze starter die graag nieuwe uitdagingen aangaat. Ik vind het belangrijk om zelfstandig verantwoordelijkheid te nemen, maar werk ook graag samen met anderen om tot betere oplossingen te komen. De combinatie van vakinhoudelijke ontwikkeling, trainingen, coaching en samenwerking met andere Young Professionals maakt dit programma voor mij een aantrekkelijke start van mijn carrière.\\
\vspace{1em}
Graag zet ik mijn technische kennis, leervermogen en enthousiasme in om mij binnen Defensie verder te ontwikkelen en bij te dragen aan betrouwbare en innovatieve IT. Ik licht mijn motivatie en achtergrond graag verder toe in een persoonlijk gesprek.\\
\vspace{1em}
Met vriendelijke groet,\\
Lucas Furer
	};
	%\node[anchor=north east, align=left] 
	%	at ($(eduEntryA.north west) + (\horGapLarge,0)$)
	%{
	%	\textbf{2023-2026}
	%};



	\coordinate (eduEnd) at ($(eduStart |- eduEntryA.south)$);
	\draw[background2, line width=1pt]
		($(eduStart)$) --
		($(eduEnd)$);

	




















\end{tikzpicture}

\end{document}
